%%%%%%%%%%%%%%%%%%%%%%%%%%%%%%%%%%%%%%%%%%%%%%%%%%%%%%%%%%%%%
\subsubsection{本问分析与建模动机}

% 第1段：任务、输入、输出和问题类型。
问题一要求……，输入为……，需要输出……，属于……问题。

% 第2段：真正的建模困难，避免只写“问题复杂、数据很多”。
本问的困难在于……之间相互耦合／……无法直接观测／……具有严格主次关系。

% 第3段：从问题结构自然引出模型，而不是先报算法名称。
考虑到……具有……性质，本文将……表示为……，并把……与……分开处理。
因此建立……模型，其作用是……。

% 仅当多阶段关系确实不容易用文字说明时加入流程图。
% 流程图应表示数学状态、决策和信息流，不使用固定框数，也不画
% “读取数据—运行程序—保存结果”一类无建模信息的工程流程。

%%%%%%%%%%%%%%%%%%%%%%%%%%%%%%%%%%%%%%%%%%%%%%%%%%%%%%%%%%%%%
\subsubsection{数据与模型准备}

% 数据处理放在首次使用相关数据的位置；没有数据清洗时删除本段。
原始数据包括……。经检查，发现……。本文采用……处理，其理由是……。
处理前后样本量／变量单位／关键统计量为……。

% 模型或算法选择须有实验、理论性质或文献依据。
% 解析、枚举、证明或题面唯一指定的方法不必强行安排算法竞赛。
候选方法包括……。根据……约束、……规模以及……验证结果，本文选择……。
与其他候选相比，该方法能够……；其限制是……。

% 可选：存在真实可比数据时使用下表。不要填写无依据的“优、良、中”。
% \begin{table}[H]
%   \centering
%   \caption{候选方法的证据化比较}
%   \begin{tabularx}{\textwidth}{CLLL}
%     \toprule
%     方法 & 约束处理或适用性 & 验证结果 & 选择结论 \\
%     \midrule
%     …… & …… & …… & …… \\
%     \bottomrule
%   \end{tabularx}
%   \label{tab:method-selection-q1}
% \end{table}

综上，本问将……作为决策对象，以……为目标，并在……约束下求解。
